% !TeX root = ../main.tex

\section{Methodology}\label{sec:methodology}

\subsection{Preliminary}

\todo{Define membership inference attack (MIA) for code LLMs. Define syntax conventions.}

\begin{figure*}[!t]
    \centering
    \includegraphics[width=0.98\textwidth]{fig/overview.pdf}
    \caption{Overview of \tool}
    \label{fig:overview}
\end{figure*}

\subsection{Multi-Language Syntax Conventions}

\todo{Describe how syntax conventions are collected for each target language (C, Java, JavaScript).}
\todo{Categorize conventions: data model, expressions, simple statements, compound statements.}

\begin{itemize}
    \item \textit{Data Model:} \todo{...}
    \item \textit{Expressions:} \todo{...}
    \item \textit{Simple Statements:} \todo{...}
    \item \textit{Compound Statements:} \todo{...}
\end{itemize}

\subsection{Syntax-Pruned MIA}

\todo{Phase 1: Token preprocessing.}

\todo{Phase 2: Syntax-convention-based token pruning.}

\begin{equation}
\ell : \mathcal{X} \rightarrow \{0,1\}
\end{equation}

\todo{Phase 3: Member probability calculation.}

\begin{equation}
\text{SPP}(x)=\frac{1}{|\mathcal{X}_{1}|}\sum_{i \in \mathcal{X}_{1}} -\log p(x_i)
\end{equation}
